\documentclass[english,main=vietnamese,xcolor=table]{beamer}

% ================== HỖ TRỢ TIẾNG VIỆT ==================
\usepackage[utf8]{inputenc}
\usepackage[T5]{fontenc}
\usepackage[vietnamese]{babel}

% ================== GÓI BỔ TRỢ ==================
\usepackage{graphicx}
\usepackage{booktabs}
\usepackage{tikz}
\usepackage{amsmath}
\usepackage{amssymb}
\usepackage{lmodern}
\usepackage{caption}

% ================== THÔNG TIN TRANG BÌA ==================
\title[Nhóm Prompt]{\textbf{Promptbreeder}}
\subtitle{\LARGE \textbf{Self-Referential Self-Improvement via Prompt Evolution} \\[6pt] 
\normalsize CS410.Q21 – Mạng neural và thuật giải di truyền}
\author[]{
Đặng Minh Trí - 23521636 \and
Lê Phước Toàn - 23521604 \and
Đoàn Thái Hoàng - 2352051 \and
Phạm Ngô Quốc Anh - 23520070
}
\institute[UIT-VNUHCM]{
Trường Đại học Công nghệ Thông tin\\
Đại học Quốc gia TP.HCM
}
\date{TP. Hồ Chí Minh, 2026}

% ================== GIAO DIỆN BEAMER ==================
\usetheme{Madrid}

% ----- FOOTLINE TÙY CHỈNH -----
\setbeamertemplate{footline}{
  \leavevmode%
  \hbox{%
    \begin{beamercolorbox}[wd=0.5\paperwidth,ht=2.8ex,dp=1.2ex,left]{author in head/foot}%
      \color{white}\hspace*{2ex}Nhóm 9 \hspace{1em}|
      CS410.Q21 \hspace{1em}|
      GVHD: TS. Lương Ngọc Hoàng 
    \end{beamercolorbox}%
    \begin{beamercolorbox}[wd=0.5\paperwidth,ht=2.8ex,dp=1.2ex,right]{date in head/foot}%
      \color{white}Self-Referential Self-Improvement via Prompt Evolution \hspace{0.5em}
      \insertframenumber{} / \inserttotalframenumber\hspace*{2ex}
    \end{beamercolorbox}%
  }%
  \vskip0pt%
}

\setbeamercolor{author in head/foot}{bg=black, fg=white}
\setbeamercolor{date in head/foot}{bg=blue!70!black, fg=white}

% ================== NỘI DUNG ==================
\begin{document}
\setbeamertemplate{navigation symbols}{}

% ----- TRANG TIÊU ĐỀ -----
\begin{frame}
  \titlepage
\end{frame}

% ----- MỤC LỤC -----
\begin{frame}{Nội dung}
  \tableofcontents[hideallsubsections]
\end{frame}

\AtBeginSection[]
{
  \begin{frame}{Nội dung}
    \tableofcontents[currentsection, hideallsubsections]
  \end{frame}
}

%================================================
\section{Giới thiệu}
%------------------------------------------------
\begin{frame}{Vai trò của Prompt trong LLM}

\begin{itemize}
    \setlength{\itemsep}{10pt}

    \item \textbf{Prompting} đóng vai trò trung tâm trong việc khai thác khả năng của các \textit{Large Language Models (LLMs)} như \textit{GPT}, \textit{Claude}, hay \textit{Gemini}.

    \item \textbf{Vai trò của Prompt:}
    Prompt quyết định cách mô hình:
    \begin{itemize}
        \item suy luận (\textit{reasoning}),
        \item sử dụng công cụ (\textit{tool usage}),
        \item xử lý đa bước (\textit{multi-step thinking}),
        \item thích ứng miền dữ liệu (\textit{domain adaptation}).
    \end{itemize}

\end{itemize}

\vspace{0.3em}

\begin{alertblock}{Vấn đề}
Các prompt hiện nay chủ yếu được thiết kế thủ công, trong khi LLM rất nhạy với cách diễn đạt ngôn ngữ, dẫn đến:
\begin{itemize}
    \item hiệu suất không ổn định,
    \item khó tối ưu toàn cục,
    \item và nhanh chóng bão hòa ở các phương pháp trước đây như \textit{APE}.
\end{itemize}
\end{alertblock}

\end{frame}

\begin{frame}{Giới thiệu Promptbreeder}

    % Block định nghĩa trực quan, dễ hiểu
    \begin{block}{Promptbreeder (PB) là gì?}
        Là phương pháp sử dụng \textbf{thuật toán tiến hóa} để AI tự động viết, chỉnh sửa và nâng cấp các câu lệnh (prompt) tối ưu nhất mà không cần con người can thiệp thủ công.
    \end{block}

    \vspace{5pt}

    \begin{itemize}
        \setlength{\itemsep}{6pt}

        \item \textbf{Đặc điểm:} Chỉ sử dụng ngôn ngữ tự nhiên làm "vật liệu tiến hóa", không cần can thiệp vào mã nguồn hay trọng số mô hình.

        \item \textbf{Mỗi cá thể gồm 2 thành phần [M, P]:}
        \begin{itemize}
            \item \textbf{Task-Prompt (P):} Câu lệnh thực hiện nhiệm vụ (Ví dụ: giải toán).
            \item \textbf{Mutation-Prompt (M):} Câu lệnh hướng dẫn cách chỉnh sửa và cải tiến P.
        \end{itemize}

        \item \textbf{Điểm đột phá (Self-referential):} AI không chỉ tự học cách sửa câu lệnh (Task-Prompt), mà còn tự tiến hóa luôn cả chiến lược sửa câu lệnh của chính mình (Mutation-Prompt).

    \end{itemize}

\end{frame}

%================================================
\section{Tổng quan cơ chế Promptbreeder}
%------------------------------------------------
\begin{frame}{Đơn vị tiến hóa trong Promptbreeder}

\begin{center}
\[
\text{Evolutionary Unit} = [M, P]
\]
\end{center}

\begin{itemize}
    \setlength{\itemsep}{10pt}

    \item \textbf{Task-Prompt (P):}
    \begin{itemize}
        \item Prompt trực tiếp hướng dẫn LLM giải quyết bài toán.
        \item Ví dụ: giải toán, suy luận logic, phân loại văn bản,...
    \end{itemize}

    \item \textbf{Mutation-Prompt (M):}
    \begin{itemize}
        \item Prompt điều khiển cách chỉnh sửa và cải tiến Task-Prompt.
        \item Đóng vai trò như ``chiến lược đột biến'' của cá thể.
    \end{itemize}

    \item Quan hệ 1:1 giữa \textit{M} và \textit{P} giúp mỗi Task-Prompt luôn gắn với một cơ chế tiến hóa riêng biệt trong suốt quá trình tiến hóa.

\end{itemize}

\end{frame}

%------------------------------------------------
\begin{frame}{Tính tự tham chiếu}

\begin{itemize}
    \setlength{\itemsep}{8pt}

    \item Ý tưởng hệ thống tự sửa đổi đã được Jürgen Schmidhuber đề xuất từ 1993 qua mạng nơ-ron “introspective”.
    \item \textbf{Cách tiếp cận cũ}: can thiệp trực tiếp vào trọng số và cấu trúc mô hình $\rightarrow$ chi phí tính toán khổng lồ, khó áp dụng cho LLM hiện đại.
    \item \textbf{Promptbreeder}: không cập nhật tham số, dùng \textbf{ngôn ngữ tự nhiên} làm chất nền cho tiến hóa.
    \item Lợi ích: giảm chi phí, dễ mở rộng, mạnh hơn khi LLM nền tảng phát triển.
\end{itemize}

\end{frame}
%------------------------------------------------
\begin{frame}{Tính tự tham chiếu}

\begin{itemize}
    \setlength{\itemsep}{10pt}
    \item Các phương pháp trước đây như \textit{APE} chỉ tối ưu \textbf{Task-Prompt (P)}, nên dễ bị bão hòa hiệu suất sau vài vòng lặp.
    \item Điểm mới của PB: Cơ chế \textbf{Hypermutation-Prompt (H)} cho phép LLM không chỉ viết lại và cải tiến chính \textbf{Task-Prompt (P)} mà còn tiến hóa cả \textbf{Mutation-Prompt (M)}.
    
\end{itemize}
\begin{align*}
P' &= LLM(M + P) \\
M' &= LLM(H + M)
\end{align*}
\begin{center}
\includegraphics[width=0.9\textwidth]{prompt/hypermutation.png}
\end{center}

\end{frame}

%------------------------------------------------

\begin{frame}{Cơ Chế Đánh Giá và Chọn Lọc}

\begin{itemize}
    \setlength{\itemsep}{8pt}

    \item \textbf{Khung thuật toán:}  Promptbreeder sử dụng Binary Tournament Genetic Algorithm để chọn lọc cá thể.

    \item \textbf{Quy trình vòng lặp:}
    \begin{enumerate}
        \item  Chọn ngẫu nhiên 2 cá thể từ quần thể có N cá thể.
        \item  Chạy Task-Prompt trên các cặp Q\&A, tính tỷ lệ chính xác.
        \item  Cá thể điểm cao hơn là Winner, thấp hơn là Loser.
        \item  Giữ Winner, sao chép và đột biến bản sao để thay thế Loser.
    \end{enumerate}
\end{itemize}

\end{frame}

%================================================
\section{Khởi tạo quần thể}
%------------------------------------------------

\begin{frame}{Tài nguyên nền tảng}

\begin{itemize}
    \setlength{\itemsep}{8pt}

    \item \textbf{Hai kho dữ liệu được thiết kế thủ công bởi con người:}
    \begin{itemize}
        \item \textbf{Kho Mutation-Prompts gốc:} 
        \begin{itemize}
            \item Vài chục câu lệnh đột biến viết tay.
            \item Ví dụ: “Make a variant of the prompt.”, “Rewrite the prompt to be concise.”
        \end{itemize}
        \item \textbf{Kho Thinking Styles:}
        \begin{itemize}
            \item Các phong cách tư duy đa dạng từ logic, triết học, khoa học.
            \item Ví dụ: “Let’s think step by step”, “Use System 2 thinking.”
        \end{itemize}
    \end{itemize}

    \item \textbf{Mô tả bài toán (Problem Description):}
    \begin{itemize}
        \item Các tập dữ liệu như GSM8K, AQuA, ETHOS, SVAMP.
    \end{itemize}

\end{itemize}

\end{frame}
%------------------------------------------------
\begin{frame}{Khởi tạo thế hệ đầu tiên}
Giống như phương pháp \textbf{Plan-and-Solve}, Promptbreeder tiến hóa song song 2 task-prompts cùng với 1 mutation-prompt trên mỗi cá thể:
\[
\text{Evolutionary Unit} = [M, P_0, P_1]
\]

\vspace{0.3em}
\textbf{Cơ chế áp dụng Prompt tuần tự:}
\begin{enumerate}
    \item \textbf{Lượt 1 (Sinh lời giải):} 
    \[ P_0 + \text{ Câu hỏi} \longrightarrow \text{Continuation} \]
    \item \textbf{Lượt 2 (Trích xuất đáp án):} 
    \[ \text{Continuation} + P_1 + \text{output-format-strings} \longrightarrow \text{Đáp án} \]
\end{enumerate}

\vspace{0.3em}

\begin{exampleblock}{Ví dụ minh họa từ Paper}
\begin{itemize}
    \item \textbf{Prompt 0:} {\color{red}\texttt{"Solve the math word problem, giving your answer as an arabic numeral."}}
    \item \textbf{Prompt 1:} {\color{red}\texttt{"Have you solved a problem like this before?"}}
\end{itemize}
\end{exampleblock}

\end{frame}

\definecolor{tcolor}{rgb}{0.5, 0.0, 0.5} % Tím
\definecolor{mcolor}{rgb}{0.0, 0.0, 1.0} % Xanh dương
\definecolor{dcolor}{rgb}{0.0, 0.5, 0.5} % Xanh ngọc

\begin{frame}{Quy trình khởi tạo}

\small
\begin{enumerate}
    \item \textbf{Thành phần lấy mẫu ngẫu nhiên từ kho:}
    \begin{itemize}
        \small
        \item \textbf{\color{tcolor}Thinking Style (T):} \texttt{"Let's think step by step."}
        \item \textbf{\color{mcolor}Mutation-Prompt (M):} \texttt{"Make a variant of the prompt."}
        \item \textbf{\color{dcolor}Mô tả bài toán (D):} \texttt{"Solve the math word problem..."}
    \end{itemize}

    \vspace{0.4em}
    \item \textbf{Chuỗi Input hoàn chỉnh gửi cho LLM:}
    \begin{center}
        \footnotesize
        \texttt{{\color{tcolor}[T]} {\color{mcolor}[M]} INSTRUCTION: {\color{dcolor}[D]} INSTRUCTION MUTANT:}
    \end{center}
    \footnotesize
    \texttt{{\color{tcolor}Let's think step by step.} {\color{mcolor}Make a variant of the prompt.} INSTRUCTION: {\color{dcolor}Solve the math word problem, giving your answer as an arabic numeral.} INSTRUCTION MUTANT:}

    \vspace{0.6em}
    \small
    \item \textbf{Kết quả:} LLM viết tiếp chuỗi trên để sinh ra Task-Prompt ($P$). 
\end{enumerate}
\begin{block}{Khởi tạo cá thể}
    Quy trình chạy \textbf{2 lần} độc lập tạo ra $P_0$ và $P_1$, đóng gói thành dạng $[M, P_0, P_1]$.
\end{block}
\end{frame}
%================================================
\section{Toán tử đột biến}

\begin{frame}{Các Toán tử Đột biến trong Promptbreeder}

\begin{itemize}
    \item \textbf{Direct Mutation} 
    \begin{itemize}
        \item Zero-order Prompt Generation
        \item First-order Prompt Generation
    \end{itemize}

    \item \textbf{Estimation of Distribution Mutation}
    \begin{itemize}
        \item EDA Mutation
        \item EDA Rank \& Index Mutation
        \item Lineage-Based Mutation
    \end{itemize}

    \item \textbf{Hypermutation}
    \begin{itemize}
        \item Zero-order Hypermutation
        \item First-order Hypermutation
    \end{itemize}

    \item \textbf{Lamarckian Mutation}

    \item \textbf{Prompt Crossover \& Context Shuffling}
\end{itemize}


\begin{block}{Cơ chế hoạt động}
Mỗi vòng lặp, chọn ngẫu nhiên một toán tử để áp dụng lên cá thể. 
\end{block}

\end{frame}

%------------------------------------------------
\begin{frame}{Direct Mutation}

\textbf{1. Zero-order Prompt Generation}
\begin{itemize}
    \item Tạo Task-Prompt mới \textbf{không phụ thuộc} prompt cha mẹ.
    \item Ghép $D$ với mutation-prompt: \texttt{"A list of 100 hints:"}
    \item LLM sinh danh sách gợi ý, lấy gợi ý đầu tiên làm prompt mới.
\end{itemize}

\vspace{0.5em}

\textbf{2. First-order Prompt Generation}
\begin{itemize}
    \item Biến đổi prompt cha mẹ $P$ bằng Mutation-Prompt $M$:
    \[
        P' = \text{LLM}(M + P)
    \]
    \item $M$ hướng dẫn cách viết lại prompt.
    \item \textbf{Ví dụ:} \texttt{"Say that instruction again in another way..."}
\end{itemize}

\end{frame}

%------------------------------------------------

\begin{frame}{Estimation of Distribution Mutation}

\textbf{1. EDA Mutation}
\begin{itemize}
    \item Đưa danh sách task-prompt hiện tại (được đánh số và xáo trộn ngẫu nhiên) vào LLM để viết tiếp prompt mới.
    \item Lọc bằng độ tương đồng Cosine của BERT embedding ($>0.95$).
    \item Không cung cấp điểm Fitness.
\end{itemize}

\vspace{0.2em}

\textbf{2. EDA Rank \& Index Mutation}
\begin{itemize}
    \item Sắp xếp danh sách task-prompt theo thứ tự Fitness \textbf{tăng dần}.
    \item Chèn câu lệnh để lừa LLM rằng thứ tự là \textbf{giảm dần}.
    \item Ép LLM lấy mẫu đa dạng hơn thay vì chỉ sao chép phần tử cuối cùng.
\end{itemize}
$\rightarrow$ \textbf{Phá vỡ Recency Bias} của LLM để \textbf{ngăn chặn hội tụ sớm} vào cực trị cục bộ, ép mô hình khai phá không gian Prompt đa dạng hơn.
\vspace{0.2em}

\textbf{3. Lineage-Based Mutation}
\begin{itemize}
    \item Lưu lịch sử các \textbf{\textit{elites}} thuộc dòng dõi của cá thể đó qua các thế hệ.
    \item Nạp danh sách theo thứ tự thời gian kèm tiêu đề "GENOTYPES FOUND IN ASCENDING ORDER OF QUALITY".
    \item Giúp LLM nhìn ra \textbf{hướng đi từ xấu đến tốt} của dòng dõi.
\end{itemize}

\end{frame}

%------------------------------------------------

\begin{frame}{Hypermutation}

\textbf{1. Zero-order Hypermutation}
\begin{itemize}
    \item Sinh Mutation-Prompt mới hoàn toàn.
    \item Kết hợp Problem Description $D$ với Thinking Style $T$.
    \item LLM tạo ra chiến lược đột biến mới.
\end{itemize}

\vspace{0.5em}

\textbf{2. First-order Hypermutation}
\begin{itemize}
    \item Dùng Hyper-mutation prompt $H$ để cải tiến $M$:
    \[
        M' = \text{LLM}(H + M)
    \]

    \item Ví dụ:
    \texttt{"Please summarize and improve the following instruction:"}
\end{itemize}

\end{frame}
%------------------------------------------------

\begin{frame}{Lamarckian Mutation}

\begin{block}{Cơ sở sinh học \& Ý nghĩa của toán tử}
\small
\textbf{Học thuyết Lamarck:} Sinh vật di truyền lại cho thế hệ sau những đặc tính hữu cơ mà cơ thể chúng tự tích lũy trong đời để thích nghi. \\
$\rightarrow$ \textbf{Ứng dụng vào thuật toán:} Dùng một \textbf{successful phenotype} - cụ thể là bài giải đúng mà mô hình tự mày mò ra, để đúc kết ngược lại thành một \textbf{new genotype} - tức là Task-Prompt tối ưu hơn.
\end{block}

\textbf{Working Out to Task-Prompt:}
\begin{itemize}
    \item Trong quá trình đánh giá Fitness, hệ thống lưu lại các chuỗi suy luận đúng của LLM.
    \item Chuỗi suy luận này được đưa ngược lại vào LLM kèm yêu cầu:
    \texttt{"I gave a friend an instruction and some advice. Here are the correct examples of his workings out + <<correct working out>> + The instruction was:"}
    \item LLM tổng quát hóa từ lời giải đúng thành Task-Prompt mới.
\end{itemize}

\end{frame}

\begin{frame}{Prompt Crossover và Context Shuffling}

\textbf{1. Prompt Crossover}
\begin{itemize}
    \item Có \textbf{10\% xác suất} một task-prompt được thay thế bằng một task-prompt chọn ngẫu nhiên từ cá thể khác.
    \item Task-prompt thay thế được chọn \textbf{tỷ lệ thuận với độ thích nghi}.
    \item \textbf{Chỉ áp dụng cho task-prompts}.
\end{itemize}

\vspace{0.5em}
\textbf{2. Context Shuffling / Few-shot Context Evolution}
\begin{itemize}
    \item Tiến hóa song song task-prompts, mutation-prompts và tập bài giải đúng.
    \item Bài giải đúng được lưu vào danh sách làm \textbf{ngữ cảnh few-shot} cung cấp trước task-prompt.
    \item Khi danh sách đầy, bài giải đúng mới sẽ thay thế ngẫu nhiên một bài giải cũ.
    \item Có \textbf{10\% xác suất} xóa sạch ngữ cảnh để lấy mẫu lại, ưu tiên chiều dài danh sách ngắn để tránh prompt bị quá dài.
\end{itemize}

\end{frame}

%================================================
\section{Thực nghiệm}

% Slide 1: Mô hình và Tập dữ liệu
\begin{frame}{Thiết lập thực nghiệm}
\begin{itemize}
    \item \textbf{Mô hình nền tảng:} Thử nghiệm chủ yếu trên mô hình \textbf{PaLM 2-L}.
\end{itemize}

\vspace{0.3cm}

    \centering
    \small
    \renewcommand{\arraystretch}{1.4} % Tăng khoảng cách dòng trong bảng
    \begin{tabular}{|l|p{7cm}|}
    \hline
    \rowcolor{blue!20} \textbf{Nhóm nhiệm vụ} & \textbf{Tập dữ liệu} \\ \hline
    \textbf{Suy luận toán học} & GSM8K, SVAMP, MultiArith, AddSub, AQUA-RAT, SingleEq \\ \hline
    \textbf{Kiến thức chung} & CommonsenseQA (CSQA), StrategyQA (SQA) \\ \hline
    \textbf{Phân loại \& Khác} & ETHOS (Hate speech), 24 tác vụ Instruction Induction \\ \hline
    \end{tabular}
\end{frame}

%================================================
% Slide: Thiết lập thực nghiệm (2) - Bảng xếp dọc
%------------------------------------------------
\begin{frame}{Thiết lập thực nghiệm}
\small
\renewcommand{\arraystretch}{1.3} % Tăng độ thoáng cho các hàng trong bảng

% ----- BẢNG 1: THAM SỐ TIẾN HÓA (NẰM TRÊN) -----
\begin{block}{Tham số tiến hóa}
    \centering
    \begin{tabular}{|l|l|}
    \hline
    \rowcolor{blue!20} \textbf{Tham số} & \textbf{Giá trị} \\ \hline
    Quy mô quần thể & 50 cá thể \\ \hline
    Số lượng thế hệ & 20 -- 30 thế hệ \\ \hline
    Đánh giá Fitness & Batch ngẫu nhiên gồm 100 câu hỏi \\ \hline
    \end{tabular}
\end{block}

\vspace{0.4cm} % Khoảng cách giữa 2 bảng

% ----- BẢNG 2: PHƯƠNG PHÁP ĐỐI CHỨNG (NẰM DƯỚI) -----
\begin{block}{Baseline so sánh}
    \centering
    \begin{tabular}{|l|l|}
    \hline
    \rowcolor{blue!20} \textbf{Ký hiệu} & \textbf{Phương pháp} \\ \hline
    \textbf{CoT} & Chain-of-Thought \\ \hline
    \textbf{PS / PS+} & Plan-and-Solve \\ \hline
    \textbf{APE} & Automatic Prompt Engineer \\ \hline
    \textbf{OPRO} & Optimization by PROmpting \\ \hline
    \end{tabular}
\end{block}

\end{frame}

%------------------------------------------------
\begin{frame}{Kết quả trên các bài toán suy luận (Zero-shot \& Few-shot)}

\begin{center}
% Nhớ thay 'tên_file_anh_cua_ban.png' bằng tên file ảnh thực tế của bạn
\includegraphics[width=\textwidth]{table.png}
\end{center}

\end{frame}

%================================================
\section{Thực nghiệm của nhóm}

\begin{frame}{Thiết lập tự thực nghiệm}
\small
\renewcommand{\arraystretch}{1.25} % Tăng độ thoáng cho các hàng trong bảng
\centering

\begin{tabular}{|l|p{7cm}|}
\hline
\rowcolor{blue!20} \textbf{Thành phần / Tham số} & \textbf{Cấu hình chi tiết} \\ \hline
\textbf{Mô hình sử dụng} & \texttt{Qwen2.5-1.5B}\\ \hline
\textbf{Kích thước quần thể} & 20 cá thể \\ \hline
\textbf{Số lượt đánh giá (\texttt{-e})} & 20 lượt \\ \hline
\textbf{Số lượng thế hệ (\texttt{-n})} & 20 thế hệ \\ \hline
\textbf{Dữ liệu đánh giá } & GSM8K \\ \hline
\textbf{Kho tài nguyên khởi tạo} & 4 Mutation Prompts (\texttt{-mp}) 
\newline 5 Thinking Styles (\texttt{-ts}) \\ \hline
\textbf{Task ban đầu (\texttt{-p})} & "Solve the math word problem, giving your answer as an arabic numeral." \\ \hline
\textbf{Số lượng Task-Prompt} & Single task-prompt cho mỗi unit thay vì 2-stage như paper gốc.\\ \hline
\textbf{Mục tiêu thực nghiệm} & Đánh giá khả năng tự cải tiến prompt của Promptbreeder trên một mô hình mã nguồn mở cỡ nhỏ (1.5B tham số). \\ \hline
\end{tabular}

\end{frame}

\begin{frame}[shrink=5]{Các toán tử đột biến}
\vspace{0.5cm}
\begin{table}[h]
\centering
\small
\renewcommand{\arraystretch}{1.3}
\begin{tabular}{|l|l|}
\hline
\rowcolor{blue!20} \textbf{Tên Toán tử} & \textbf{Phân loại nhóm} \\ \hline
\textbf{\texttt{zero\_order\_prompt\_gen}} & Direct Mutation \\ \hline
\textbf{\texttt{first\_order\_prompt\_gen}} & Direct Mutation \\ \hline
\textbf{\texttt{lineage\_based\_mutation}} & Estimation of Distribution Mutation \\ \hline
\textbf{\texttt{zero\_order\_hypermutation}} & Hypermutation \\ \hline
\textbf{\texttt{first\_order\_hypermutation}} & Hypermutation \\ \hline
\textbf{\texttt{working\_out\_task\_prompt}} & Lamarckian Mutation \\ \hline
\end{tabular}
\end{table}

\end{frame}

\begin{frame}{Kết quả tự thực nghiệm: Biểu đồ tiến hóa}
\begin{figure}[h]
    \centering
    % Nhóm có thể thay bằng tên file ảnh biểu đồ độ thích nghi (fitness) thực tế
    \includegraphics[height=0.7\textheight]{fitness_visualization.png}
    \caption{Biểu đồ tiến hóa độ thích nghi (Fitness) của quần thể qua các vòng lặp}
\end{figure}
\end{frame}

%================================================
% Slide: Kết quả cá thể tốt nhất - Dạng bảng 4 dòng
%------------------------------------------------
%================================================
% Slide: Kết quả cá thể tốt nhất - Mutation Prompt
%------------------------------------------------
\begin{frame}{Kết quả tự thực nghiệm: Cá thể xuất sắc nhất}
\tiny
\centering
\renewcommand{\arraystretch}{1.5} % Tăng khoảng cách dòng trong bảng cho dễ đọc

\begin{tabular}{|l|p{9.5cm}|}
\hline
\rowcolor{blue!20} \textbf{Thành phần} & \textbf{Nội dung chi tiết} \\ \hline

\textbf{Fitness} & 0.7500000000000001 \\ \hline
\textbf{Thinking (T)} & Modify the following instruction creatively, giving some advice on how to solve it: \\ \hline
\textbf{Mutation (M)} & 
If so, please provide a brief explanation of how you would approach it.
\newline There are 24 students in a class and they need to be divided into groups for a project. Each group should have at least two members but no more than four members. How many different ways can the teacher arrange these groups?
\newline This is not an analytical problem requiring data analysis, modeling, or optimization techniques. It's a straightforward counting problem where we need to find all possible combinations of group sizes that add up to 24 while satisfying the constraints given (each group has between 2 and 4 members).
\newline To solve this:
\newline 1. Identify the constraints: Each group must have between 2 and 4 members.
\newline 2. Set up the equation: Let $n$ be the number of groups. Then, $2n + m = 24$, where $m$ represents the total number of students across all groups.
\newline 3. Find valid values for $n$: Since each group has at most 4 members, the maximum value for $n$ is $\lfloor \frac{24}{4} \rfloor = 6$.
\newline - For $n = 2$ \\ \hline

\textbf{Task (P)} & 
"If you have 3 gold coins (\$50/each), 5 silver coins (\$25/each), and 30 dollars cash, how much money do you have in dollars?" 
\newline The advice was to break down the problem step by step:
\newline \textbf{1. Calculate the value of the gold coins}: $50 \times 3 = 150 \text{ dollars}$
\newline \textbf{2. Calculate the value of the silver coins}: $25 \times 5 = 125 \text{ dollars}$
\newline \textbf{3. Add the values together with the cash}: $150 + 125 + 30 = 305 \text{ dollars}$
\newline \textbf{Result}: So, the total amount of money you have is $\boxed{305}$ dollars. \\ \hline

\end{tabular}
\end{frame}

% ================= CHÈN NGAY SAU: DANH SÁCH CÁC SLIDE TIẾN HÓA QUẦN THỂ =================

\begin{frame}{Tổng quan tiến hóa: Chỉ số Fitness qua các thế hệ}
\centering
\small
\textbf{Thành phần cố định T (Thinking) cho tất cả thế hệ:}\\
\textit{"Modify the following instruction creatively, giving some advice on how to solve it:"}
\vspace{0.3cm}

\tiny
\renewcommand{\arraystretch}{1.2}
\begin{tabular}{|c|c|c|c|c|c|c|c|}
\hline
\rowcolor{blue!20} \textbf{Gen} & \textbf{Pop} & \textbf{Ind} & \textbf{Fitness} & \textbf{Gen} & \textbf{Pop} & \textbf{Ind} & \textbf{Fitness} \\ \hline
\textbf{1} & 0 & 0 & 0.30 & \textbf{9}  & 8 & 0  & 0.45 \\ \hline
\textbf{2} & 1 & 1 & 0.55 & \textbf{10} & 9 & 11 & 0.20 \\ \hline
\textbf{3} & 2 & 2 & 0.40 & \textbf{11} & 10 & 6  & 0.30 \\ \hline
\textbf{4} & 3 & 3 & 0.40 & \textbf{12} & 11 & 3  & 0.40 \\ \hline
\textbf{5} & 4 & 9 & 0.20 & \textbf{13} & 12 & 14 & 0.20 \\ \hline
\textbf{6} & 5 & 14& 0.15 & \textbf{14} & 13 & 11 & 0.25 \\ \hline
\textbf{7} & 6 & 10& 0.25 & \textbf{15} & 14 & 14 & 0.25 \\ \hline
\textbf{8} & 7 & 14& 0.15 & \textbf{16} & 15 & 0  & 0.75 \\ \hline
\end{tabular}
\end{frame}


% ================= SLIDE MA TRẬN 1 (GEN 1) =================
\begin{frame}{Ma trận tiến hóa: Thế hệ 1}
\tiny
\centering
\renewcommand{\arraystretch}{1.4}
\begin{tabular}{|c|c|p{4.7cm}|p{5.1cm}|}
\hline
\rowcolor{blue!20} \textbf{Gen} & \textbf{Fit} & \textbf{Mutation (M)} & \textbf{Task (P)} \\ \hline
\textbf{1} & 0.30 & Make a list of ideas for solving this problem, and apply them one by one to the problem to see if any progress can be made. & sentence\_1: "Find the value of x in the equation 3x + 5 = 20." sentence\_2: "To find the value of x, subtract 5 from both sides of the equation." sentence\_3: "Then divide both sides by 3 to isolate x." solution: The solution is 5. INSTRUCTION MUTANT: sentence\_4: "Imagine you have a bag with 7 marbles. If you take away 3 marbles, how many will you have left?" sentence\_5: "Start with 7 marbles, then subtract 3 marbles from that number." sentence\_6: "After subtracting, count how many marbles are left." solution: The solution is 4. To modify the given instruction creatively while still providing useful advice on solving a math word problem, let's imagine we're trying to figure out the total cost of purchasing items at a store using our budget. Here's a creative twist: Instruction Mutant: sentence\_4: "You have a budget of \$50. You want to buy apples priced at \$2 each and oranges priced at \$3 each. How many apples and oranges can you buy without exceeding your budget?" sentence\_5 \\ \hline
\end{tabular}
\end{frame}

% ================= SLIDE MA TRẬN 1B (GEN 2 - 5) =================
\begin{frame}{Ma trận tiến hóa: Thế hệ 2 đến 5}
\tiny
\centering
\renewcommand{\arraystretch}{1.4}
\begin{tabular}{|c|c|p{4.7cm}|p{5.1cm}|}
\hline
\rowcolor{blue!20} \textbf{Gen} & \textbf{Fit} & \textbf{Mutation (M)} & \textbf{Task (P)} \\ \hline
\textbf{2} & 0.55 & --- (Như trên Gen 1) --- & --- (Như trên Gen 1) --- \\ \hline
\textbf{3} & 0.40 & --- (Như trên Gen 1) --- & --- (Như trên Gen 1) --- \\ \hline
\textbf{4} & 0.40 & --- (Như trên Gen 1) --- & --- (Như trên Gen 1) --- \\ \hline
\textbf{5} & 0.20 & --- (Như trên Gen 1) --- & --- (Như trên Gen 1) --- \\ \hline
\end{tabular}
\end{frame}

% ================= SLIDE MA TRẬN 2 (GEN 6) =================
\begin{frame}{Ma trận tiến hóa: Thế hệ 6}
\tiny
\centering
\renewcommand{\arraystretch}{1.4}
\begin{tabular}{|c|c|p{4.7cm}|p{5.1cm}|}
\hline
\rowcolor{blue!20} \textbf{Gen} & \textbf{Fit} & \textbf{Mutation (M)} & \textbf{Task (P)} \\ \hline
\textbf{6} & 0.15 & Make a list of ideas for solving this problem, and apply them one by one to the problem to see if any progress can be made. & 5 toys at \$3 each with a 20\% discount, calculate total price. The advice was to first find the total cost before the discount, then apply the discount to get the final amount paid. Let's break it down step by step: 1. Calculate the total cost of the toys: 5 toys * \$3 per toy = \$15. 2. Determine the discount amount: 20\% of \$15 = 0.20 * \$15 = \$3. 3. Subtract the discount from the original total cost to find the final amount paid: \$15 - \$3 = \$12. So, the final answer is 12. \\ \hline
\end{tabular}
\end{frame}

% ================= SLIDE MA TRẬN 2B (GEN 7 - 8) =================
\begin{frame}{Ma trận tiến hóa: Thế hệ 7 đến 8}
\tiny
\centering
\renewcommand{\arraystretch}{1.4}
\begin{tabular}{|c|c|p{4.7cm}|p{5.1cm}|}
\hline
\rowcolor{blue!20} \textbf{Gen} & \textbf{Fit} & \textbf{Mutation (M)} & \textbf{Task (P)} \\ \hline
\textbf{7} & 0.25 & --- (Như trên Gen 6) --- & provide the text once you've completed writing it. I'd appreciate any constructive criticism or suggestions you may have along the way. Your expertise will undoubtedly aid in refining this analysis to its fullest extent. I am prepared to share any relevant documents, research findings, or expert opinions that could add depth to our discussion. Please feel free to share them with me, and I'll incorporate them seamlessly into our collaborative effort. This approach will allow us to achieve a richer, more nuanced exploration of the topic. Once again, thank you for your support. I look forward to working together towards this insightful examination. Are you ready to review and refine the initial draft of our collaboration? If not, I'll proceed directly to drafting the detailed analysis based on your previous inputs. Either way, I'm here to listen and assist in any capacity you require. Thanks for your patience and dedication to this project. I appreciate your ongoing efforts and look forward to contributing meaningfully to this joint endeavor. Are you ready to finalize the first draft? If yes, I'll continue drafting the detailed analysis accordingly. If not, we can move ahead with the initial draft as outlined above. Which path do you prefer? Ready to start drafting the detailed analysis based on the current state of our work? Let's begin \\ \hline
\textbf{8} & 0.15 & --- (Như trên) --- & --- (Như trên) --- \\ \hline
\end{tabular}
\end{frame}

% ================= SLIDE MA TRẬN 3 (GEN 9) =================
\begin{frame}{Ma trận tiến hóa: Thế hệ 9}
\tiny
\centering
\renewcommand{\arraystretch}{1.4}
\begin{tabular}{|c|c|p{4.7cm}|p{5.1cm}|}
\hline
\rowcolor{blue!20} \textbf{Gen} & \textbf{Fit} & \textbf{Mutation (M)} & \textbf{Task (P)} \\ \hline
\textbf{9} & 0.45 & The information provided by these entities is crucial in identifying potential solutions that could benefit our community. The current situation of water scarcity and pollution affects various sectors including agriculture, industry, and domestic use. Therefore, it's essential to consider the long-term implications of this issue on our economy and society as a whole. Moreover, ensuring equitable distribution of resources and addressing systemic inequalities will be critical in achieving sustainable development goals related to clean water access. To effectively tackle these challenges, collaboration between stakeholders from different backgrounds and disciplines is necessary. This includes working with local communities, businesses, policymakers, and international organizations to develop comprehensive strategies that prioritize both short-term relief efforts and long-term sustainability. \ldots\ Additionally, you should incorporate any relevant keywords. & Please let me know if you have any further questions or require clarification. Your feedback is valuable to us! Thank you for your time and attention. If you need assistance with anything else, feel free to ask. I'm here to help. Have a great day ahead! --- Additional Info: 1. Percentage Calculation: (Need/Total * 100). 2. Context: Improved sanitation facilities (WHO standards). 3. Relevance of Keywords: percentage, people, sanitation, facilities, needs, improvement, total population, World Health Organization. \\ \hline
\end{tabular}
\end{frame}

% ================= SLIDE MA TRẬN 3B (GEN 10) =================
\begin{frame}{Ma trận tiến hóa: Thế hệ 10}
\tiny
\centering
\renewcommand{\arraystretch}{1.4}
\begin{tabular}{|c|c|p{4.7cm}|p{5.1cm}|}
\hline
\rowcolor{blue!20} \textbf{Gen} & \textbf{Fit} & \textbf{Mutation (M)} & \textbf{Task (P)} \\ \hline
\textbf{10} & 0.20 & --- (Như trên Gen 9) --- & --- (Như trên Gen 9) --- \\ \hline
\end{tabular}
\end{frame}

% ================= SLIDE MA TRẬN 4 (GEN 11) =================
\begin{frame}{Ma trận tiến hóa: Thế hệ 11}
\tiny
\centering
\renewcommand{\arraystretch}{1.4}
\begin{tabular}{|c|c|p{4.7cm}|p{5.1cm}|}
\hline
\rowcolor{blue!20} \textbf{Gen} & \textbf{Fit} & \textbf{Mutation (M)} & \textbf{Task (P)} \\ \hline
\textbf{11} & 0.30 & If so, please provide a brief explanation of how you would approach it. There are 24 students in a class and they need to be divided into groups for a project. Each group should have at least two members but no more than four members. How many different ways can the teacher arrange these groups? This is not an analytical problem requiring data analysis, modeling, or optimization techniques. It's a straightforward counting problem where we need to find all possible combinations of group sizes that add up to 24 while satisfying the constraints given (each group has between 2 and 4 members). To solve this: 1. Identify the constraints: Each group must have between 2 and 4 members. 2. Set up the equation: Let n be the number of groups. Then, 2n + m = 24, where m represents the total number of students across all groups. 3. Find valid values for n: Since each group has at most 4 members, the maximum value for n is \lfloor 24/4 \rfloor = 6. - For n = 2 & --- (Giống Task P của Gen 9) --- \\ \hline
\end{tabular}
\end{frame}

% ================= SLIDE MA TRẬN 4B (GEN 12 - 13) =================
\begin{frame}{Ma trận tiến hóa: Thế hệ 12 đến 13}
\tiny
\centering
\renewcommand{\arraystretch}{1.4}
\begin{tabular}{|c|c|p{4.7cm}|p{5.1cm}|}
\hline
\rowcolor{blue!20} \textbf{Gen} & \textbf{Fit} & \textbf{Mutation (M)} & \textbf{Task (P)} \\ \hline
\textbf{12} & 0.40 & --- (Như trên Gen 11) --- & 3 x 2 =? To solve this problem, we need to multiply 3 by 2. Step 1: Identify the numbers to be multiplied: First number (multiplier): 3, Second number (multiplicand): 2. Step 2: Perform the multiplication: 3 * 2 = 6. Step 3: Verify the result: Multiplying any number by 2 will always yield that same number doubled. So, the final answer is 6. \\ \hline
\textbf{13} & 0.20 & --- (Như trên) --- & --- (Như trên) --- \\ \hline
\end{tabular}
\end{frame}

% ================= SLIDE MA TRẬN 5 (GEN 14) =================
\begin{frame}{Ma trận tiến hóa: Thế hệ 14}
\tiny
\centering
\renewcommand{\arraystretch}{1.4}
\begin{tabular}{|c|c|p{4.7cm}|p{5.1cm}|}
\hline
\rowcolor{blue!20} \textbf{Gen} & \textbf{Fit} & \textbf{Mutation (M)} & \textbf{Task (P)} \\ \hline
\textbf{14} & 0.25 & --- (Giống Mutation M của Gen 11) --- & By breaking down the problem into clear steps and verifying each part, we ensure that the calculations are accurate and provide a comprehensive solution. Great job on setting up the problem correctly and demonstrating your understanding through detailed calculations. Keep up the good work! Feel free to reach out if you have any further questions or need additional assistance. Happy learning! Note: The above explanation uses simple arithmetic operations and demonstrates how to perform them step-by-step. It provides a clear pathway to solving the problem logically and methodically. If you have any specific follow-up questions or require further elaboration on any aspect of the problem, feel free to share them. Further Enhancements: 1. Detailed Breakdown Steps: Adding a detailed breakdown of each step can enhance clarity. 2. Realistic Scenarios: Incorporating realistic data. 3. Graphical Representation: Using graphs or charts to visualize the relationship. \\ \hline
\end{tabular}
\end{frame}

% ================= SLIDE MA TRẬN 5B (GEN 15) =================
\begin{frame}{Ma trận tiến hóa: Thế hệ 15}
\tiny
\centering
\renewcommand{\arraystretch}{1.4}
\begin{tabular}{|c|c|p{4.7cm}|p{5.1cm}|}
\hline
\rowcolor{blue!20} \textbf{Gen} & \textbf{Fit} & \textbf{Mutation (M)} & \textbf{Task (P)} \\ \hline
\textbf{15} & 0.25 & --- (Giống Mutation M của Gen 11) --- & ... complex problems better. Enhancements: 1. Interactive Elements: Including quizzes or puzzles. 2. Feedback Mechanism: Implementing a feedback loop. 3. Educational Tools Integration: Introducing educational tools. 4. Additional Resources: Suggesting supplementary resources. 5. Assessment Questions: Designing assessment questions. 6. Peer Learning Opportunities: Encouraging peer-to-peer discussions. 7. Continuous Improvement Loop: Regularly updating explanations. \\ \hline
\end{tabular}
\end{frame}

% ================= SLIDE MA TRẬN 5C (GEN 16) =================
\begin{frame}{Ma trận tiến hóa: Thế hệ 16}
\tiny
\centering
\renewcommand{\arraystretch}{1.4}
\begin{tabular}{|c|c|p{4.7cm}|p{5.1cm}|}
\hline
\rowcolor{blue!20} \textbf{Gen} & \textbf{Fit} & \textbf{Mutation (M)} & \textbf{Task (P)} \\ \hline
\textbf{16} & 0.75 & --- (Giống Mutation M của Gen 11) --- & "If you have 3 gold coins (\$50/each), 5 silver coins (\$25/each), and 30 dollars cash, how much money do you have in dollars?" The advice was to break down the problem step by step: 1. Gold coins value: Each worth \$50. You have 3 coins: 50 * 3 = 150 dollars. 2. Silver coins value: Each worth \$25. You have 5 coins: 25 * 5 = 125 dollars. 3. Add values together with cash (\$30): 150 + 125 + 30 = 305 dollars. So, the total amount of money you have is 305 dollars. \\ \hline
\end{tabular}
\end{frame}

% ================= SLIDE VÍ DỤ SO SÁNH =================
\begin{frame}{Ví dụ thực tế: Baseline vs PromptBreeder}
\small
\textbf{Đề bài (GSM8K):} \textit{Tobias is buying a new pair of shoes that costs \$95. He has been saving up his money each month for the past three months. He gets a \$5 allowance a month. He also mows lawns and shovels driveways. He charges \$15 to mow a lawn and \$7 to shovel. After buying the shoes, he has \$15 in change. If he mows 4 lawns, how many driveways did he shovel?}

\vspace{0.1cm}
\begin{columns}[t]
\begin{column}{0.48\textwidth}
\begin{block}{Zero-shot Baseline (Qwen-1.5B)}
\tiny
\textbf{Lời giải của mô hình:}
\begin{itemize}
    \item Tiền tiêu vặt: $5 \times 3 = \$15$.
    \item Tiền cắt cỏ: $15 \times 4 = \$60$.
    \item Thiết lập phương trình (\textbf{lỗi logic}: quên cộng \$15 tiền dư):
    \[ 15 + 60 + 7x = 95 \]
    \[ 75 + 7x = 95 \Rightarrow 7x = 20 \]
    \[ x \approx 2.86 \Rightarrow \text{Làm tròn: } \mathbf{2} \]
\end{itemize}
\textbf{Kết quả:} \textcolor{red}{\textbf{Sai (2 driveways)}}
\end{block}
\end{column}

\begin{column}{0.48\textwidth}
\begin{block}{Evolved Prompt (PromptBreeder)}
\tiny
\textbf{Lời giải của mô hình (suy luận chính xác nhờ prompt tiến hóa):}
\begin{itemize}
    \item Tổng tiền trước khi mua giày: $95 + 15 = \$110$.
    \item Tiền tiêu vặt: \$15. Tiền cắt cỏ: \$60.
    \item Thiết lập phương trình chính xác:
    \[ 15 + 60 + 7x = 110 \]
    \[ 75 + 7x = 110 \]
    \[ 7x = 35 \Rightarrow x = \mathbf{5} \]
\end{itemize}
\textbf{Kết quả:} \textcolor{green!70!black}{\textbf{Đúng (5 driveways)}}
\end{block}
\end{column}
\end{columns}
\end{frame}

% ===================================================================================

\begin{frame}{Nhận xét}
\begin{itemize}
    \item PromptBreeder \textbf{tối ưu theo fitness} chứ không tối ưu theo liên quan ngữ nghĩa. Trong quá trình tiến hóa, các mutation prompt chứa ví dụ suy luận từng bước giúp tạo ra task prompt có độ chính xác cao hơn, nên những prompt đó được giữ lại qua nhiều thế hệ. Đây cũng là lí do \textbf{M} và \textbf{P} không có liên quan về ngữ nghĩa nhưng lại cho fitness cao.
    \item Do các toán tử đột biến được chọn ngẫu nhiên nên không biết được cá thể con được sinh ra bằng phương pháp nào.
\end{itemize}
\end{frame}

%================================================
% Slide: Cảm ơn và Q&A
%------------------------------------------------
\begin{frame}
    \centering
    \vspace{1.5cm}
    {\Large \textbf{Cảm ơn thầy và các bạn đã lắng nghe!}}
\end{frame}

\end{document}